%%%%%%%%%% Arquivo de configuração do projeto
%% Insira seus dados e modifique o modelo com suas preferências
%% Comente/descomente linhas para modificar o modelo
%% Aviso! O modelo não segue uma norma específica

%%%%%%%%%%%%%% Informações dos projeto
\tipo{Plano de Trabalho}
\titulo{Otimização Combinatória na Geração de Árvores de Decisão}
\aluno{Sofia Valverde Villas Bôas}
\orientador{Prof. Dr. Rafael Crivellari Saliba Schouery}
\palavraschave{Estudo Dirigido, Otimização Combinatória, Machine Learning}


%%%%%%%%%%%%%% Pacotes extras
\usepackage{lipsum} % Gera texto dummy
\usepackage{fullpage} % Uniformiza as margens (1,5 cm) 
\usepackage{hyperref}

%%%%%%%%%%%%%% Bibliografia
%% Para referências estilo ABNT tire o comentário de uma das linhas 
%% Manual do pacote de referências: https://www.abntex.net.br/

%\usepackage[num]{abntex2cite}	% Referências numéricas padrão ABNT
%\usepackage[alf]{abntex2cite}  % Referências alfabéticas padrão ABNT

\AtEndDocument{
    \bibliografia{bibtex.bib} % Arquivo com o bibtex
    \estilobibliografico{plain} % Estilo das referências se não usar ABNT
%% Mais estilos: https://www.overleaf.com/learn/latex/bibtex_bibliography_styles
}

%% Para adicionar backrefs tire o comentário da linha abaixo,
%% backrefs listam as paginas em que uma referência foi citada
%\usepackage[brazilian,hyperpageref]{backref} % Adiciona backrefs

%%%%%%%%%%%%%% Confs
\AtBeginDocument{
    \maketitle % Gera cabeçalho
    \pagestyle{plain} % Efetivamente, insere a numeração
    \renewcommand{\equationautorefname}{Equação}
    \renewcommand{\figureautorefname}{Figura}
    \renewcommand{\tableautorefname}{Tabela}
    \renewcommand{\sectionautorefname}{Seção}
    \renewcommand{\subsectionautorefname}{Subseção}
    \renewcommand{\theoremautorefname}{Teorema}
}


%%%%%%%%%%%%%%%% Comandos úteis para matemática
\newcommand{\Z}{\mathbb{Z}}
\newcommand{\Q}{\mathbb{Q}}
\newcommand{\R}{\mathbb{R}}
\newcommand{\trans}{^\intercal}

\DeclareMathOperator{\posto}{posto}
\DeclareMathOperator{\capacidade}{cap}
\DeclareMathOperator{\valor}{val}
\DeclareMathOperator{\mdc}{mdc}
\DeclareMathOperator{\adj}{adj}
\DeclareMathOperator{\conv}{conv}
\DeclareMathOperator{\cone}{cone}

\DeclarePairedDelimiter\ceil{\lceil}{\rceil}
\DeclarePairedDelimiter\floor{\lfloor}{\rfloor}

\newtheorem{teorema}{Teorema}
\newtheorem{definicao}{Definição}
\newtheorem{lema}{Lema}
\newtheorem{proposicao}{Proposição}
\newtheorem{corolario}{Corolário}
\newtheorem{fato}{Fato}
\theoremstyle{definition}